\section{State Machine and Mission Logic}\label{sec:statemachine}

The mission controller is implemented as a deterministic finite state machine (FSM)~\cite{wagner2006fsm}, a widely adopted pattern in autonomous robotics for its verifiability, traceability, and ease of runtime monitoring. The FSM orchestrates the entire flight --- from power-on through search, detection, operator verification, payload delivery, and return --- while enforcing safety invariants at every transition.

\subsection{Software Architecture}\label{subsec:arch}

The mission software separates concerns across seven Python modules, each with a single responsibility (see also Section~\ref{sec:safety} for the safety implications of this decomposition):

\begin{itemize}
    \item \texttt{config.py} --- every tunable parameter (altitudes, speeds, camera, NFZ polygons, connection strings);
    \item \texttt{states.py} --- a \texttt{State} class defining 20 string constants, one per discrete state;
    \item \texttt{vision.py} --- camera acquisition and AI detection behind a single interface: \texttt{detect\_in\_image(frame)} $\rightarrow$ \texttt{(found, x, y, conf)};
    \item \texttt{planning.py} --- lawnmower search pattern generator, returns an ordered list of GPS waypoints;
    \item \texttt{navigation.py} --- MAVLink command abstraction (\texttt{goto}, \texttt{set\_speed}, \texttt{send\_velocity}, \texttt{land}, \texttt{arm});
    \item \texttt{geofence.py} --- NFZ boundary distance computation and repulsive offset generation (Section~\ref{sec:geofencing});
    \item \texttt{state\_machine.py} --- per-state handler methods, packaged as a mixin class.
\end{itemize}

\noindent The orchestration lives in \texttt{main.py}, which defines the central class \texttt{VisualFlightMission}. This class inherits from \texttt{StateHandlersMixin} (defined in \texttt{state\_machine.py}) using Python's mixin pattern: the mixin contributes all per-state handler methods (e.g.\ \texttt{\_handle\_search}, \texttt{\_handle\_verify}) and shared helpers, while the main class owns the event loop, telemetry, HUD rendering, and module wiring. The mixin declares no \texttt{\_\_init\_\_} and relies on attributes initialised by the consuming class, keeping the two files decoupled yet composable --- the orchestrator at approximately 820~lines, the handler mixin at approximately 940~lines.

\paragraph{Dispatch Table.}
Rather than a chain of \texttt{if}/\texttt{elif} branches, the main loop uses a dictionary that maps each \texttt{State} constant to its handler method:

\begin{verbatim}
_dispatch = {
    State.INIT:              self._handle_init,
    State.SEARCH:            self._handle_search,
    State.CENTERING:         self._handle_centering,
    State.VERIFY:            self._handle_verify,
    ...                      # 18 entries total
}

handler = _dispatch.get(self.state)
if handler:
    handler(target_found, px_u, px_v, key)
\end{verbatim}

\noindent This provides $O(1)$ dispatch, eliminates fall-through errors, and makes it straightforward to add or remove states without modifying the dispatch logic.

\subsection{State Definitions}\label{subsec:states}

The system defines 20~discrete states, grouped into four functional phases: startup, search, engagement, and recovery. Two states (\textsc{Manual} and \textsc{Done}) are handled outside the dispatch table as they require special control flow.

\paragraph{Startup Phase.}
\begin{itemize}
    \item \textbf{INIT} --- Entry state. Establishes a MAVLink connection to the autopilot via the configured transport (serial on the Raspberry Pi, UDP for SITL). Transitions to \textsc{Connecting} once the connection object is created.
    \item \textbf{CONNECTING} --- Waits for a valid heartbeat from the Cube autopilot and requests data streams at 10\,Hz. A 15\,s warning is issued if no heartbeat arrives. Transitions to \textsc{Arming} on first heartbeat.
    \item \textbf{ARMING} --- Sets GUIDED flight mode, waits for a GPS fix ($\geq$\,3D, $\geq$\,6~satellites), records the home position, and sends the arm command. A 120\,s timeout alerts the operator to check pre-arm conditions. Transitions to \textsc{Takeoff} once motors are armed.
    \item \textbf{TAKEOFF} --- Sends \texttt{MAV\_CMD\_NAV\_TAKEOFF} to the configured search altitude (\texttt{TARGET\_ALT}\,=\,35\,m). Monitors altitude and re-enters \textsc{Arming} if the vehicle disarms unexpectedly (a known SITL artefact) or transitions to \textsc{Done} on real hardware. Transitions to \textsc{Pre\_Waypoints} or \textsc{Transit\_To\_Search} when 90\% of target altitude is reached.
\end{itemize}

\paragraph{Search Phase.}
\begin{itemize}
    \item \textbf{PRE\_WAYPOINTS} --- Flies a sequence of operator-defined transit waypoints (loaded from JSON) at transit speed before entering the search area. Provides a controlled ingress path.
    \item \textbf{TRANSIT\_TO\_SEARCH} --- Navigates to the first lawnmower waypoint at transit speed (\texttt{TRANSIT\_SPEED\_MPS}\,=\,15\,m/s). Transitions to \textsc{Search} on arrival ($<$\,2\,m).
    \item \textbf{SEARCH} --- Executes the lawnmower pattern waypoint-by-waypoint. Each frame is processed by the vision system; confirmed detections are filtered against rejected targets, no-fly zones, and the search boundary, then enqueued. When a valid queued detection exists, the aircraft departs the pattern and transitions to \textsc{Centering}. If all waypoints are exhausted without a confirmed target, a rescan pass at lower altitude may be triggered (configurable), or the mission ends.
\end{itemize}

\paragraph{Engagement Phase.}
\begin{itemize}
    \item \textbf{CENTERING} --- Navigates to the estimated GPS position of the detected target. A lock radius (\texttt{DETECT\_LOCK\_RADIUS\_M}\,=\,5\,m) prevents the aircraft from diverting to secondary detections that appear during approach; new detections outside this radius are queued for later investigation. A 60\,s timeout returns the aircraft to \textsc{Search} if convergence fails. Transitions to \textsc{Verify} when within 1\,m of the target.
    \item \textbf{DESCENDING} --- Reserved for future altitude-staged verification. Currently passes through directly to \textsc{Verify}.
    \item \textbf{VERIFY} --- The operator-in-the-loop decision point (Section~\ref{subsec:operator}). The aircraft holds position above the target while the operator classifies it as confirmed (\texttt{Y}), rejected (\texttt{N}), or item of interest (\texttt{I}). A 120\,s timeout auto-rejects the target to prevent indefinite hovering.
    \item \textbf{APPROACH} --- Flies to the computed landing spot, offset 7.5\,m from the target in an operator-selected cardinal direction (N/E/S/W), at 3\,m altitude. Transitions to \textsc{Hover\_Target} on arrival.
    \item \textbf{HOVER\_TARGET} --- Executes a 15\,s payload deployment sequence using a servo-actuated release mechanism (two-stage: partial release at $t=3$\,s, full release at $t=6$\,s). The servo is returned to its closed position before departure. Transitions to \textsc{Return\_Transit} or \textsc{Return\_Home}.
\end{itemize}

\paragraph{Recovery Phase.}
\begin{itemize}
    \item \textbf{RETURN\_TO\_SEARCH} --- Returns to the point where the aircraft departed the search pattern after a rejected or interest-classified target, then resumes the lawnmower pattern.
    \item \textbf{RETURN\_FROM\_MANUAL} --- Returns to the position where manual override was engaged, then resumes the previously active state.
    \item \textbf{RETURN\_TRANSIT} --- Retraces the pre-waypoint transit path in reverse order at search altitude. Provides controlled egress from the search area.
    \item \textbf{RETURN\_HOME} --- Navigates to the recorded home (take-off) position. A safety guard verifies that a GPS fix was successfully obtained during arming; if not, the aircraft lands in place rather than flying to a potentially erroneous location.
    \item \textbf{LANDING} --- Sends \texttt{MAV\_CMD\_NAV\_LAND} and monitors descent via three independent touchdown mechanisms (Section~\ref{subsec:safety}). Transitions to \textsc{Done}.
    \item \textbf{DONE} --- Terminal state. Logs the landing error (distance from home) and mission elapsed time.
    \item \textbf{HOVER} --- Fallback loiter state entered if no search waypoints are available after take-off, with a 60\,s timeout leading to \textsc{Done}.
    \item \textbf{MANUAL} --- Software manual override, accessible from any active state (Section~\ref{subsec:safety}).
\end{itemize}

\subsection{State Transitions}\label{subsec:transitions}

% Placeholder for state machine diagram
\begin{figure}[htbp]
    \centering
    % \includegraphics[width=\textwidth]{figures/state_machine_diagram.pdf}
    \fbox{\parbox{0.9\textwidth}{\centering\vspace{3cm}\textit{State machine diagram placeholder}\vspace{3cm}}}
    \caption{State transition diagram for the SAR mission controller. The primary mission path flows left-to-right across the top; engagement branches downward from \textsc{Search}. Manual override (dashed) is accessible from any active state. Safety transitions (dotted) can interrupt from any airborne state.}
    \label{fig:state_machine}
\end{figure}

The main loop executes at approximately 50\,Hz and follows a fixed five-stage pipeline each iteration: (1)~read telemetry from the autopilot, (2)~run the vision pipeline and render the HUD, (3)~process operator key/button input, (4)~invoke the dispatch handler for the current state, and (5)~enforce geofence constraints. This ordering ensures that geofence enforcement always runs \emph{after} the state handler, so any waypoint command issued in step~4 can be overridden by a repulsive velocity in step~5. Table~\ref{tab:transitions} summarises the principal transitions.

\begin{table}[htbp]
\centering
\caption{State transition table. Conditions marked with $^*$ are safety-critical.}
\label{tab:transitions}
\small
\begin{tabular}{lll}
\toprule
\textbf{From} & \textbf{Condition} & \textbf{To} \\
\midrule
INIT & MAVLink connection created & CONNECTING \\
CONNECTING & First heartbeat received & ARMING \\
ARMING & GPS fix $\geq$3D, $\geq$6 sats, motors armed & TAKEOFF \\
TAKEOFF & Altitude $\geq$ 90\% of target & PRE\_WAYPOINTS \\
PRE\_WAYPOINTS & All transit waypoints reached & TRANSIT\_TO\_SEARCH \\
TRANSIT\_TO\_SEARCH & Within 2\,m of first search waypoint & SEARCH \\
SEARCH & Valid detection dequeued & CENTERING \\
SEARCH & All waypoints exhausted & DONE / rescan \\
CENTERING & Within 1\,m of target GPS & VERIFY \\
CENTERING & 60\,s timeout$^*$ & SEARCH \\
VERIFY & Operator presses Y + side selected & APPROACH \\
VERIFY & Operator presses N & RETURN\_TO\_SEARCH \\
VERIFY & Operator presses I & RETURN\_TO\_SEARCH \\
VERIFY & 120\,s timeout$^*$ (auto-reject) & SEARCH \\
APPROACH & Within 2\,m of landing spot, alt $<$ 4\,m & HOVER\_TARGET \\
HOVER\_TARGET & 15\,s elapsed (payload deployed) & RETURN\_TRANSIT \\
RETURN\_TRANSIT & All transit waypoints retraced & RETURN\_HOME \\
RETURN\_HOME & Within 2\,m of home position & LANDING \\
LANDING & Touchdown detected or timeout$^*$ & DONE \\
\midrule
Any active state & M key pressed & MANUAL \\
MANUAL & M key pressed & RETURN\_FROM\_MANUAL \\
RETURN\_FROM\_MANUAL & Within 3\,m of departure point & \textit{previous state} \\
Any active state$^*$ & GPS degraded $>$ 5\,s & LANDING (via RTL) \\
Any active state$^*$ & MAVLink link lost & DONE (via RTL) \\
Any active state$^*$ & Inside NFZ boundary & MANUAL \\
\bottomrule
\end{tabular}
\end{table}

Every transition is logged with the old state, new state, and mission elapsed time in the format \texttt{STATE: OLD --> NEW [T+MM:SS]}. All state-specific timeout flags are reset on entry (via the centralised \texttt{\_set\_state} method), preventing stale timeout warnings from persisting across state re-entries.

\subsection{Detection Queue}\label{subsec:queue}

During the search phase, multiple targets may be detected before any can be investigated. The system manages these using a FIFO detection queue with three properties:

\begin{enumerate}
    \item \textbf{Spatial deduplication} --- Before enqueuing, each detection is checked against rejected targets, items of interest, and existing queue entries using a configurable radius (\texttt{REJECTED\_TARGET\_RADIUS\_M}, default 5\,m). Detections within this radius of any known position are silently discarded, preventing the queue from filling with repeated sightings of the same object.
    \item \textbf{Bounded capacity} --- The queue is capped at \texttt{MAX\_DETECT\_QUEUE} entries (default 20). When full, the oldest entry is dropped and logged. This prevents unbounded memory growth during long flights over cluttered terrain.
    \item \textbf{Validation on dequeue} --- When a target is popped for investigation, it is re-validated against the NFZ polygon, search area boundary, and rejection list. Targets that have become invalid since enqueuing (e.g.\ a nearby target was subsequently rejected) are skipped. This lazy validation avoids the overhead of re-checking the entire queue every iteration.
\end{enumerate}

An optional \texttt{--smart-detect} mode requires \texttt{DETECT\_CONFIRM\_FRAMES} consecutive positive frames (default~3) before enqueuing, suppressing transient false positives at the cost of slightly increased detection latency. During manual override, detections continue to be queued rather than discarded, so returning to autonomous mode can immediately investigate targets observed while the operator was flying manually.

\subsection{Operator-in-the-Loop Verification}\label{subsec:operator}

The \textsc{Verify} state implements a human-in-the-loop decision gate, motivated by the high cost of false positives in SAR operations~\cite{murphy2014disaster}. Autonomous descent onto an incorrectly identified target wastes battery endurance, delays area coverage, and risks damage to the aircraft or bystanders. The operator classifies each detection into one of four categories:

\begin{itemize}
    \item \textbf{Confirm (Y)} --- The target is the casualty. When \texttt{--center-verify} mode is active, GPS averaging is performed over 10 seconds to refine the position estimate using an inverse-variance weighted mean of all samples collected while hovering. The operator then selects a landing side (N/E/S/W) to avoid overflying the casualty during approach.
    \item \textbf{Reject (N)} --- The target is a false positive. Its GPS position is added to a rejection list; future detections within a configurable radius are suppressed, preventing repeated investigation of the same object.
    \item \textbf{Interest (I)} --- The target is noteworthy but not the primary casualty (e.g., debris, equipment). Its position is logged for later review while the search continues.
    \item \textbf{False Positive (X)} --- Equivalent to reject, used in the web dashboard for ergonomic clarity.
\end{itemize}

A 120\,s timeout ensures the aircraft does not hover indefinitely if the operator is unreachable; the target is automatically rejected and the search resumes. Escalating warnings are issued at 60, 90, and 100\,s, with a live countdown displayed on the HUD. This tiered warning structure provides the operator with progressively urgent cues:

\begin{enumerate}
    \item At 60\,s: first warning (``60s remaining'');
    \item At 90\,s: second warning (``30s remaining'');
    \item At 100\,s onward: warnings every 10 seconds until timeout.
\end{enumerate}

The operator interface supports three input modalities: keyboard over SSH/PuTTY, browser buttons on the web dashboard (Section~\ref{sec:groundstation}), and OpenCV key events for simulation. All three feed into a unified command queue, ensuring identical behaviour regardless of input source. This supervisory control paradigm~\cite{goodrich2007human} lets autonomy handle the high-bandwidth search while human judgement is reserved for the high-consequence classification decision.

\subsection{Safety Overrides and Fault Handling}\label{subsec:safety}

Five independent safety mechanisms operate across all flight states, forming a defence-in-depth architecture~\cite{ardupilot_failsafe}.

\paragraph{Manual Override (M key).}
Pressing M at any time transitions the FSM to \textsc{Manual}, immediately halting autonomous commands by sending a zero-velocity command. The departure position and altitude are recorded so that pressing M again triggers \textsc{Return\_From\_Manual}, which navigates back to the departure point before resuming the previous state. If the operator observes a valid detection during manual flight, it is queued; upon exiting manual mode, the system checks for queued detections and transitions directly to \textsc{Centering} if one is available, or returns to the departure point if none exist. WASD keys provide \texttt{MANUAL\_FLY\_SPEED\_MPS}\,=\,5\,m/s velocity commands, R/F control vertical movement at \texttt{MANUAL\_CLIMB\_RATE\_MPS}\,=\,2\,m/s, and Q/E rotate the yaw. This affords the operator complete manual authority without leaving the software framework.

\paragraph{RC Kill Switch.}
The RC transmitter can override the flight controller at the hardware level, independent of the companion computer. Switching the RC mode channel to \texttt{STABILIZE} or \texttt{LOITER} immediately removes the aircraft from GUIDED mode. The software monitors the \texttt{HEARTBEAT.custom\_mode} field and warns when the mode is no longer GUIDED, but takes no corrective action --- the safety pilot retains full authority. This layered approach ensures that a software fault cannot prevent the safety pilot from regaining control~\cite{ardupilot_failsafe}.

\paragraph{Return-to-Launch on Link Loss.}
Two link-loss scenarios are handled. First, if the MAVLink \texttt{recv\_match} call raises a connection exception (socket error or broken pipe), an emergency RTL command (mode~6) is sent as a best-effort measure and the mission terminates. Second, GPS degradation is continuously monitored: if the fix type drops below 3D or the satellite count falls below six for more than 5\,s, an emergency RTL is triggered and the FSM transitions to \textsc{Landing}. ArduPilot's own GCS failsafe (\texttt{FS\_GCS\_ENABLE}) provides an additional firmware-level RTL if heartbeats from the companion computer cease entirely.

\paragraph{Geofence Enforcement.}
A no-fly zone (NFZ) polygon, defined from the site survey KML file corresponding to the SSSI boundary, is enforced through three complementary layers:

\begin{enumerate}
    \item \textbf{Plan-time filtering} --- the lawnmower planner's \texttt{filter\_waypoints()} method removes any waypoint within \texttt{NFZ\_WAYPOINT\_BUFFER\_M}\,=\,30\,m of the NFZ boundary, preventing the search pattern from ever routing the aircraft into the restricted area.
    \item \textbf{Runtime speed capping} --- at every loop iteration, the aircraft's distance to the NFZ boundary is computed using \texttt{cv2.pointPolygonTest} for signed-distance evaluation. Within \texttt{NFZ\_SLOW\_ZONE\_M}\,=\,20\,m, speed is linearly reduced from \texttt{NFZ\_ZONE\_MAX\_SPEED\_MPS}\,=\,3.0\,m/s at the zone edge down to \texttt{NFZ\_MIN\_SPEED\_MPS}\,=\,0.3\,m/s at the boundary. If the aircraft crosses the hard boundary (\texttt{NFZ\_HARD\_BOUNDARY\_M}\,=\,3\,m inside the polygon), the FSM forces an immediate transition to \textsc{Manual} and issues a zero-velocity halt command.
    \item \textbf{Repulsive potential field} --- within \texttt{NFZ\_INNER\_RANGE\_M}\,=\,23\,m of an inner polygon offset 20\,m inside the NFZ boundary, a constant-magnitude repulsive velocity of \texttt{NFZ\_PUSH\_SPEED\_MPS}\,=\,3.0\,m/s is applied away from the nearest boundary segment. This deflects the aircraft from the boundary edge even when the current waypoint lies close to it. The repulsive vector is also rendered on the HUD as a visual indicator for the operator.
\end{enumerate}

\noindent Detected targets whose estimated GPS positions fall inside the NFZ or outside the search area boundary are silently discarded rather than investigated.

\paragraph{Landing Fault Tolerance.}
The \textsc{Landing} state employs three independent touchdown detection mechanisms to address the well-documented barometric drift problem in small UAS:

\begin{enumerate}
    \item \textbf{ArduPilot auto-disarm} --- the autopilot's accelerometer-based landing detector disarms the motors when sustained low vibration is detected after a land command, requiring no companion-computer involvement;
    \item \textbf{Altimeter threshold} --- if the barometric altitude remains below 0.5\,m for 5 consecutive loop ticks, the companion computer sends a disarm command directly;
    \item \textbf{Absolute timeout} --- a 90\,s wall-clock timeout triggers a forced disarm (using the MAVLink emergency disarm parameter 21196), ensuring the aircraft does not hover indefinitely if barometric drift keeps the reported altitude above the threshold.
\end{enumerate}

\noindent If the land command does not produce descent within 5\,s, it is retried up to five times before forced disarm. This layered approach ensures that landing completes under all observed failure modes, including a barometric drift deadlock documented during bench testing.

\paragraph{State Timeouts.}
Every blocking state carries a wall-clock timeout: \textsc{Arming} warns at 120\,s, \textsc{Takeoff} at 60\,s, \textsc{Centering} aborts to \textsc{Search} at 60\,s, \textsc{Verify} auto-rejects at 120\,s, and \textsc{Hover} transitions to \textsc{Done} at 60\,s. Consequently, no state can hang the mission indefinitely. The centralised \texttt{\_set\_state} method resets all timeout flags and records the entry timestamp on every transition, so each state's timeout logic is self-contained and immune to stale values from previous visits.
